\ifx\allfiles\undefined
%生成目录超链接，使用UTF8字符集
\documentclass[12pt,UTF8]{ctexart}


\usepackage{mathtools}%花括号括起的公式
\usepackage{enumerate}%列表
\usepackage{graphicx}%表格
\usepackage[table]{xcolor}%彩色表格
\usepackage{booktabs}%调整表格行高


\renewcommand\tabcolsep{50pt}
\renewcommand\arraystretch{2}

%文本框
\usepackage{tcolorbox}
\tcbuselibrary{skins,vignette,raster,listings,listingsutf8,theorems, breakable,magazine,fitting}

%页面设置
\usepackage{geometry}%设置页面
\usepackage{setspace}%设置局部行距

\usepackage{ctex}
\usepackage{CJK}

%图像处理
\usepackage{graphicx}
\usepackage{float}%允许浮动
\usepackage{caption}%设置标题
\usepackage{graphicx, subfig}

%附录
\usepackage[toc,page]{appendix}

%导言区
\title{中小微企业的信贷决策}

%不要页眉
\pagestyle{plain}
\bibliographystyle{plain}

\geometry{a4paper}%将纸张设置为A4大小

%修改文档结构层次的名称
\usepackage{titlesec}%chapter1修改为一
\titleformat{\part}{\centering\Huge\bfseries}{\thepart}{1em}{}
\renewcommand\thepart{\chinese{part}}

%subsubsubsection的使用
%\usepackage{titlesec}
%\titleclass{\subsubsubsection}{straight}[\subsection]
%
%\newcounter{subsubsubsection}[subsubsection]
%\renewcommand\thesubsubsubsection{\thesubsubsection.\arabic{subsubsubsection}}
%\renewcommand\theparagraph{\thesubsubsubsection.\arabic{paragraph}}
%
%\titleformat{\subsubsubsection}
%{\normalfont\normalsize\bfseries}{\thesubsubsubsection}{1em}{}
%\titlespacing*{\subsubsubsection}
%{0pt}{3.25ex plus 1ex minus .2ex}{1.5ex plus .2ex}
%
%\makeatletter
%\renewcommand\paragraph{\@startsection{paragraph}{5}{\z@}%
%	{3.25ex \@plus1ex \@minus.2ex}%
%	{-1em}%
%	{\normalfont\normalsize\bfseries}}
%\renewcommand\subparagraph{\@startsection{subparagraph}{6}{\parindent}%
%	{3.25ex \@plus1ex \@minus .2ex}%
%	{-1em}%
%	{\normalfont\normalsize\bfseries}}
%\def\toclevel@subsubsubsection{4}
%\def\toclevel@paragraph{5}
%\def\toclevel@paragraph{6}
%\def\l@subsubsubsection{\@dottedtocline{4}{7em}{4em}}
%\def\l@paragraph{\@dottedtocline{5}{10em}{5em}}
%\def\l@subparagraph{\@dottedtocline{6}{14em}{6em}}
%\makeatother
%
%\setcounter{secnumdepth}{4}

\begin{document}
\else
\fi
%分文件开始


\begin{abstract}

	\parskip=6pt
	%开始：题头
	综合考虑企业的实力与信誉，量化分析得到企业的信贷风险，我们就能给出银行对这些企业的信贷策略，本文拟解决相关问题。
	%结束：题头


	%开始：问题1的摘要
	对于问题一，我们用商业智能软件Power BI对附件1中的123家企业的企业信息、进项发票信息、销项发票信息进行分析，得到反映经营规模的有效销项金额和反映盈利能力的经营利润率这两项反映企业实力的指标、反映供给侧稳定性的正数有效进项发票率和反映需求侧稳定性的正数有效销项发票率这两项反映企业供求关系稳定性的指标，再将企业信誉评级与违约与否进行量化。据处理所得的指标，我们建立层次分析模型\cite{ref1}，求出上述指标在信贷风险中的权重，即可求得信贷风险。最后，我们用MATLAB构建规划模型，求得银行在年度信贷总额固定时对企业的信贷策略，结果以表格的形式呈现了不同企业在不同年度信贷总额下的信贷数额。
	%结束：问题1的摘要


	%开始：问题2的摘要	
	对于问题二，我们用Power BI对附件2中的302家企业的进项发票和销项发票信息进行数据分析，求得问题一所提到的各项指标。为评出信誉评级，我们在MPai数据科学平台上用问题一中各项指标和信誉评级来训练决策树\cite{ref2}，并用决策树评定出信誉评级。接着，用各项指标和信誉评级加权求和得到信贷风险，最后构建规划模型求得银行在年度信贷总额为1亿元时对这些企业的信贷策略，结果以表格的形式呈现了不同企业的信贷数额。
	%结束：问题2的摘要


	%开始：问题3的摘要		
	对于问题三，考虑疫情对企业的影响。我们据企业名称将302家企业划分为13个行业，并将这些行业受疫情影响程度分别设定为有利、基本无影响、不利和非常不利这四类，随后用层次分析法得出不同影响程度对信贷风险的权重，将问题二中的信贷风险加权后，得到在疫情影响下的信贷风险。最后，我们构建规划模型\cite{ref1}求得在疫情影响下银行在年度信贷总额为1亿元时的信贷调整策略，结果以表格的形式呈现了不同企业的信贷数额和调整值。
	%结束：问题3的摘要


\end{abstract}


%开始：关键词
\textbf{关键词：}
层次分析法%关键词1
\quad
数据标准化%关键词2
\quad
决策树%关键词3
\quad
多元线性规划%关键词4

%另起一段
\qquad\qquad\enspace%使用空格进行缩进
MATLAB%关键词5
\quad
商业智能%关键词6
%结束：关键词


\newpage


%分文件结束
\ifx\allfiles\undefined
\end{document}
\fi